\documentclass[12pt,a4paper]{article}

\usepackage[top=25mm,bottom=25mm,left=25mm,right=25mm]{geometry}
\usepackage{xeCJK}
\usepackage{fontspec}
\usepackage{setspace}

\setmainfont{Times New Roman}
\setCJKmainfont{IPAexMincho}

\setstretch{1.5}
\setlength{\parindent}{2em}
\pagestyle{empty}

\begin{document}

\begin{center}
{\LARGE 「お疲れ様」}\\[1em]
{\large 賈 康麗}
\end{center}

\vspace{1em}
	小さい子供のころから、おじいさん、おばあさん、父、母に囲まれて成長した。兄も姉もいるし、幸福な家族だった。でも、四五歳の私にとってそれは足りなかった。家族の役に立ちたかった。
	秋。窓から外に蕭々とふっている風に揺れている木の下に薪を割っている父の姿を見つめた。上着を取って外に出かけた。
	「お父さんを助けたい。」
	「いらないよ。あなたはまだ小さい。しかし、兄を呼んできて。」
	ときどき、私はもう少し主張すると、おじいさんは微笑みながら、
	「静かにすれば助けられるよ。」と言う。
	退屈した私は、さまよいながら無意識に台所に着いた。おばあさんは夕飯を作っているのだ。さっくり、さっくりと専門家みたいに大根を薄く切っていた。
	「おばあちゃん、お手伝いできますか。」と聞いた
	「いいよ。」削皮刀をくれた。「ニンジンの皮をむいてなさい。」
	私は熱心に三本のニンジンの皮をむいて、そして五胡のジャガイモも皮をむいた。おばあさんに見せた。
	「はい、お疲れ、」と言いながら、おばあさんは自分でジャガイモの皮をむき始めた。
	「みら、こちらにまだ皮があるよ。」
	私はため息をついた。どうして自分には能力がないの。早く大人になりたい！
	子供のころそんな感情は大人を助けたいという憧れだった。あるいはただ好奇心のせいかもしれない。年を取るにしたがって、祖父母、父、母は私の助けが本当にいるようになってきた。以前、私は家族の支援に頼っていた。今は、私が家人を保護することができる。

\end{document}